\documentclass[12pt]{article}

\usepackage[margin=1 in]{geometry}
\usepackage{amsmath, amsthm, amssymb, scrextend}
\usepackage{enumitem}
\usepackage{mathtools}
\usepackage{mdframed}
\usepackage{tikz}
\usepackage{listings}
\usepackage{bm}
\usepackage{dsfont}
\usepackage{xcolor}
\usepackage{soul}

\setlength{\parindent}{0pt}

\newenvironment{soln}
    {\begin{mdframed}}
    {\end{mdframed}}

\DeclareMathOperator{\Var}{Var}
\DeclareMathOperator{\E}{\mathbf{E}}
\DeclareMathOperator{\Cov}{Cov}
\DeclareMathOperator{\1}{\mathds{1}}

\renewcommand{\qedsymbol}{$\blacksquare$}

\begin{document}
\title{Concepts from PH 240B For This Project}
\author{Candus Shi}
\maketitle

\section{Roadmap of Statistical Learning}
Question of interest: Among midlife, menopausal hormone therapy (MHT)-naive women in SWAN who are free of hypertension at baseline, how does the cumulative incidence of hypertension by visit $K$ differ under rules for initiating MHT at different stages of the menopausal transition and response to vasomotor symptoms?

\hfill

Shorter version: How does the timing of menopausal hormone therapy (MHT) initiation affect the risk of hypertension over $K$ years of follow-up among midlife women in SWAN?
\begin{enumerate}
	\item Data: $n$ iid observations of random variable $O \sim P_0$, where 
	\[
	O = (W, L(1), A(1), C(1), Y(1), \dots, L(10), A(10), C(10), Y(10))
	\]
	where
	\begin{itemize}
		\item $W$ are baseline covariates: age, race/ethnicity, education, baseline blood pressure
		\item $L(t)$ are time-varying covariates at visit $t$: menopausal status, vasomotor symptoms, height, weight, alcohol use, smoking, physical activity, diabetes status, depression, anxiety, missingness indicator of exposure and outcome 
		\item $A(t)$ is exposure at visit $t$: use of MHT since the last visit (combination estrogen/progestin, estrogen pills, estrogen injection or patch, or progestin pills)
		\item $C(t)$ is censoring at visit $t$: first instance of right-censoring in outcome 
		\item $Y(t)$ is outcome at visit $t$: incident hypertension which is a) blood pressure $\ge$ 130/80 mmHg (\texttt{SYSBP1/DIABP1} and \texttt{SYSBP2/DIABP2}), or b) use of blood pressure/anti-hypertensive medication since visit $t$ (\texttt{BP1} and \texttt{BP2}), or c) indicator of being told by doctor of hypertension (\texttt{HIGHBP})
	\end{itemize}
	Note: the censoring variable must be created, and is only for true right-censoring, i.e. missingness in outcome until end of follow-up. 
	
	\item Statistical Model: $\mathcal M$ is nonparametric; no assumptions on the set of possible probability distributions of $O$.
	\item Target Parameter: We are interested in comparing different dynamic treatment rules for initiation of MHT.
	\begin{itemize}
		\item $d_0$: $A(t) = 0$ at every visit.
		\item $d_\text{vm}$: initiate MHT at first visit at which a woman reports vasomotor symptoms, and continue use.
		\item $d_\text{early}$: initiate MHT at first visit at which a woman is in late perimenopause or early postmenopause, and continue use.
		\item $d_\text{late}$: initiate MHT $\delta = \{1, 2, 3\}$ visits after the $d_\text{early}$ trigger, and continue use.
	\end{itemize}
	We would be interested in 
	\[
	\Psi_{d_\text{vm}}(P_0) - \Psi_{d_0}(P_0), \Psi_{d_\text{early}}(P_0) - \Psi_{d_0}(P_0), \text{ and } \Psi_{d_\text{early}}(P_0) - \Psi_{d_\text{late}}(P_0),
	\]
	where 
	\[
	\Psi_d(P_0) = \mathbf E_0[Y_{d, \overline c = 0}(K)] = P(T_d \le K)= 1 - S_d(K).
	\]
	$T_d$ and $S_d$ will be defined later.
	
	The last estimand directly addresses the timing hypothesis in MHTs which is the hypothesis that the benefits and harms of MHT initiation can depend on what stage of the menopausal transition a woman is in, with particular emphasis on the risk of cardiovascular outcomes. More simply, it is the idea that cardiovascular risk is lower if a woman takes MHT earlier in the transition whereas cardiovascular risk may be higher if a woman takes MHT later in the transition.
	\item Identification and Estimation:
	\item Inference: 
\end{enumerate}

\section{Censored data in survival analysis}
We have interval censoring, which is represented 
\[
O = (C_1, N(C_1), \dots, C_{10}, N(C_{10}))
\]
where $C_1 < \dots < C_{10}$ are monitoring times. So, if we are interested in some failure time $T$, we only know if $T$ occurs in some interval $[C_i, C_j]$. 

\section{Product integral representation of Kaplan-Meier}
\subsection{Hazard for discrete and continuous survival time}
\begin{itemize}
	\item $S_0(t) = P_0(T > t)$ is the \hl{\textbf{survival function}} at time $t$. We can plot this for $t \in [0, \dots, 10]$, and it should be a monotonically decreasing function. 
	\item $F_0(t) = P_0(T \le t)$ is the \hl{\textbf{cumulative distribution function}}. So, $S_0(t) = 1 - F_0(t)$. 
	\item If $T$ is discrete, then the \hl{\textbf{hazard}} at time $t$ is defined as 
	\[
	\lambda_0(t) = P_0(T = t \mid T \ge t) = \frac{P_0(T = t)}{S_0(t-)}
	\]
	where $f_0(t) = P_0(T = t)$ is the density of $T$ wrt the counting measure $\mu_c$: 
	\[
	P_0(T \in A) = \int_A f_0(t)d\mu_c(t) = \sum_{j: t_j \in A}f_0(t_j).
	\]
	\item If $T$ is continuous, then the \hl{\textbf{hazard}} at time $t$ is defined as
	\[
	\lambda_0(t) = \frac{f_0(t)}{S_0(t)}
	\]
	where $f_0(t)$ is the density of $T$ wrt the Lebesgue measure $\mu_\ell$:
	\[
	P_0(T \in A) = \int_A f_0(t)\,d\mu_\ell(t) = \int_A f_0(t)\,dt.
	\]
\end{itemize}
We talk about hazard because the survival function can be expressed in terms of hazard, and it is easier to estimate the hazard rather than estimate the survival directly. So, we'll estimate the hazard, then map it back to survival.

\hfill

\subsection{Cumulative hazard}
\begin{itemize}
	\item For discrete or continuous $T$, the \hl{\textbf{cumulative hazard}} is defined as 
	\[
	\Lambda_0(t) = \int_0^t \lambda_0(s)d\mu(s)
	\]
	where $\mu$ is the counting or Lebesgue measure, respectively. In our case (discrete $T$), this will be 
	\[
	\Lambda_0(t) = \sum_{j: t_j \le t}\lambda_0(t_j).
	\]
	\item More generally, we can write 
	\[
	\Lambda_0(t) = \int_0^t d\Lambda_0(s)
	\]
	where $d\Lambda_0$ is the measure generated by $\Lambda_0$. 
	\item We can also write the cumulative hazard in terms of survival:
	\[
	\Lambda_0(t) = \int_0^t \frac{dF_0(s)}{S_0(s-)}.
	\]
\end{itemize}

In plain words:
\begin{itemize}
	\item Survival at time $t$ is the probability that an individual survives past time $t$ (a woman does not have hypertension past time $t$). 
	\item Hazard at time $t$ is the intensity of death at time $t$ given that an individual survives up to time $t$ (intensity of hypertension at time $t$ given a woman does not have hypertension by time $t$). 
	\item Cumulative hazard is accumulated risk by time $t$ (TODO: elaborate).
\end{itemize}

\subsection{Product integral representation of survival in terms of cumulative hazard}
\begin{itemize}
	\item For general (discrete and continuous) survival, the survival function can written as a product integral of the cumulative hazard:
	\[
	S_0(t) = \prod_{s \in [0, t]}(1 - d\Lambda_0(s))
	\]
	\item We can think of this infinite product as a limit where we 
	\begin{enumerate}
		\item Partition $[0, t]$ into small intervals, $(a_j, b_j]$.
		\item Take the finite product of $1 - \Lambda_0((a_j, b_j])$ over all intervals.
		\item Let the width $\delta = b_j - a_j$ converge to 0. 
	\end{enumerate}
	Moreover, here is a sketch of a proof of why this works: 
	\begin{align*}
		\prod_{s \in [0, t]} (1 - d\Lambda_0(s)) &= \prod_{j = 0}^{J-1} (1 - \Lambda_0((t_j, t_{j+1}]) \\
		&= \prod_{j = 0}^{J - 1}\left(1 - \frac{P(T \in (t_j, t_{j+1}])}{P(T > t_j)}\right) \\
		&= \prod_{j = 0}^{J - 1}\frac{P(T > t_j) - P(T \in (t_j, t_{j+1}])}{P(T > t_j)} \\
		&= \prod_{j = 0}^{J - 1}\frac{P(T > t_{j+1})}{P(T > t_j)} \\
		&= \frac{P(T > t_1)}{P(T > t_0)} \cdot \frac{P(T > t_2)}{P(T > t_1)} \cdot \dots \cdot \frac{P(T > t_J)}{P(T > t_{J-1})} \\
		&= \frac{P(T > t_J)}{P(T > t_0)} \\
		&= P(T > t) \\ 
		&= S(t)
	\end{align*}
	since $P(T > t_0) = P(T > 0) = 1$. 
\end{itemize}

In the continuous case, the product integral representation of the survival function can be simplified to
\[
S_0(t) = \exp(-\Lambda_0(t)).
\]

\subsection{Kaplan-Meier estimator}
How do we estimate the survival function nonparametrically? The product integral representation of the survival function motivates the Kaplan-Meier estimator.

\hfill

For simplicity, suppose we have independent right censoring, i.e. for $n$ iid copies of 
\[
O = \big(\widetilde T = \min(T, C), \Delta = 1(T \le C)\big) \sim P_0,
\]
we assume $C \perp T$. 

\hfill

Further, let 
\begin{itemize}
	\item $G_0$ be the cdf of $C$.
	\item $\overline G_0 = 1 - G_0$ be the survival function of censoring.  
	\item $F_0$ be the cdf of $T$.
	\item $S_0 = 1 - F_0$ be the survival function.
\end{itemize}
Then, $P_0$ is implied by $F_0$ and $G_0$. So, the distribution of $(\widetilde T, \Delta)$ is parametrized by $F$ and $G$, i.e. of the form $P_{F, G}$. The joint distribution can be written as this density, 
\[
dP_{F, G}(\tilde t, \delta) = \delta dF(\tilde t) \overline G(\tilde t -) + (1 - \delta) S(\tilde t)dG(\tilde t).
\]
How to read this: since $\delta \in \{0, 1\}$, for a censored individual ($\delta = 0 \Rightarrow T > C \Rightarrow \widetilde T = C$), the joint distribution is $S(\tilde t) dG(\tilde t)$ because $S(\tilde t)$ is saying the event hasn't happened, yet, and $dG(\tilde t)$ is saying the individual got censored at time $\tilde t$. 

\hfill

Likewise, for an uncensored individual ($\delta = 1 \Rightarrow T \le C \Rightarrow \widetilde T = T$), the joint distribution is $dF(\tilde t) \overline G(\tilde t-)$ because $dF(\tilde t)$ is saying the event happened at time $\tilde t$, and $\overline G(\tilde t-)$ is saying the individual hasn't been censored, yet. 

\hfill

A helpful equivalent statement which might be useful for likelihood factorizations, is noticing that 
\[
dP_{F, G}(\tilde t, \delta) = \big[dF(\tilde t)\,\overline G(\tilde t-)\big]^{\delta}\,\big[S(\tilde t)\,dG(\tilde t)\big]^{1-\delta}.
\]

\hfill

If we let $\mathcal F$ be the set of all cdfs, then the set of possible probability distributions of $(\widetilde T, \Delta)$ can be summarized in the following statistical model,
\[
\mathcal M = \{P_{F, G}: F \in \mathcal F, G \in \mathcal F\}.
\]

Then, the full-data target parameter mapping is $\Psi^\text{full}: \mathcal F \to \mathbb R$, where
\[
\Psi^\text{full}(F) = (1 - F)(t_0),
\]
and the target parameter estimand of interest is 
\[
\psi_0^\text{full} = S_0(t_0).
\]
Using our assumption that $C \perp T$, the above target parameter can be identified by the observed data distribution. Recall that we can first estimate the cumulative hazard, then substitute it into the product integral to get the estimate for the survival.
\begin{align*}
d\Lambda_0(t) &= \frac{dF_0(t)}{S_0(t-)} \\
&= \frac{dF_0(t) \overline G_0(t-)}{S_0(t-)\overline G_0(t-)} \\
&= \frac{dP_{0, 1}(t)}{\overline P_0(t-)}
\end{align*}
where
\begin{itemize}
	\item $P_{0, 1}(t) = P_0(\widetilde T \le t, \Delta = 1) = \int_0^t dF_0(s) \overline G_0(s-)$
	\item $\overline P_0(t) = P_0(\widetilde T \ge t) = S_0(t-)\overline G_0(t-)$
\end{itemize}

Then, we can estimate the survival function via the following identification result
\[
S_0(t_0) = \prod_{s \in [0, t_0]}\left(1 - \frac{dP_{0, 1}(s)}{\overline P_0(s-)}\right).
\]

A very simple and natural estimator of the survival function estimand above is its plug-in estimator, which is also known as the \hl{\textbf{Kaplan-Meier estimator}}.
\[
S_n(t_0) = \prod_{s \in [0, t_0]}\left(1 - \frac{dP_{n, 1}(s)}{\overline P_n(s-)}\right),
\]
where
\begin{itemize}
	\item $P_{n, 1}(t) = P_n(\widetilde T \le t, \Delta = 1) = \frac1n \sum_{i=1}^n 1(\widetilde T_i \le s, \Delta_i = 1)$.
	\item $\overline P_n(t) = P_n(\widetilde T \ge t) = \frac1n \sum_{i=1}^n 1(\widetilde T_i \ge t)$.
\end{itemize}

\section{Kaplan-Meier as NPMLE}
It turns out that Kaplan-Meier is an NPMLE. 



\section{Empirical probability measure and functional CLT}

\section{Asymptotic linearity of an estimator and its influence curve}

\end{document}